\appendixsection{Параметры экспериментов}

В приложении приведены параметры окружения и запусков нагрузочного тестирования.

Кластер состоит из пяти узлов; на каждом --- два менеджера задач Apache Flink по 30
слотов (итого 10 менеджеров задач, 300 слотов); на одном узле дополнительно ---
менеджер заданий. Среда исполнения --- Java~17. Бэкенд хранения --- объектное
хранилище, совместимое с S3 (на базе Apache Ozone).
Развёртывание --- средствами Ansible.

В сценарии~1 (TPC-DS) масштаб набора --- \SI{100}{\giga\byte}; данные
загружены в Apache Paimon как таблицы только для добавления, разбиты на разделы
(типичный размер файла раздела --- порядка \SI{256}{\mega\byte}); параллелизм
запросов --- 300; на всех форматах выполнились \num{104} запроса (три прогона,
усреднение); сравнивались форматы Parquet, ORC, Vortex и Lance.

В сценарии~2 (потоковый near-real-time) параметры представительного запуска
приведены в таблице~\ref{tab:run-params}.

\begin{table}
    \caption{Параметры запуска потокового near-real-time-сценария}
    \label{tab:run-params}
    \begin{tabularx}{\textwidth}{|>{\raggedright\arraybackslash}p{0.32\textwidth}|X|}\hline
    \textbf{Параметр} & \textbf{Значение} \\ \hline
    Форматы & ORC, Parquet, Vortex \\ \hline
    Число бакетов & порядка сотни \\ \hline
    Фаза INIT & \num{50000000} строк, последовательные ключи, без ограничения скорости \\ \hline
    Фаза UPDATE & порядка \num{200000000} строк, случайные ключи, ограничение скорости $\approx 100\,000$~строк/с \\ \hline
    Интервал чекпойнтов Flink & 240~с (режим near-real-time) \\ \hline
    DURING-выборки & пакетные запросы во время записи с интервалом в несколько минут \\ \hline
    Число повторов цикла & не менее трёх \\ \hline
    Размер буфера записи & \SI{1}{\giga\byte} \\ \hline
    Размер пакета записи Vortex & \num{65536} строк \\ \hline
    Режим статистик & полные статистики по столбцам \\ \hline
    Политики хранения & расширенные (снимки 500--10000, время хранения 24~ч, асинхронная очистка) \\ \hline
    \end{tabularx}
\end{table}

Использованы следующие метрики. Для записи --- пропускная способность (строк в секунду), отдельно по
фазам INIT и UPDATE (с оговоркой о дедупликации в фазе UPDATE). Для чтения --- время
выполнения запроса в миллисекундах, отдельно по фазам DURING и FINAL, с усреднением
по выборкам и повторам. Дополнительно фиксировались число успешно завершённых
запросов под нагрузкой и наличие выбросов времени выполнения.
